% =============================================================================
% CAPÍTULO 3: PROCESSOS ORGANIZACIONAIS HABILITADORES
% Alinhado à ISO/IEC/IEEE 12207:2021 (6.2) e MPS.BR (Níveis G e F)
% Fonte: manual-processos/manual/capitulos/03-processos-organizacionais.md
% =============================================================================

\aberturacapitulo{Capítulo 3}{Processos Organizacionais Habilitadores}

\secao{Modelo de ciclo de vida adotado e justificado}

A \termochave{CineSoft Solutions} adotou um modelo de ciclo de vida \destaqueazul{híbrido e iterativo-incremental}, combinando a disciplina e o rigor de controle da \termochave{ISO/IEC/IEEE 12207:2021} com práticas ágeis baseadas no framework \textbf{Scrum/Kanban}.

Em sistemas de missão crítica como a bilheteria cinematográfica — onde bugs em produção podem ocasionar venda dupla ou paralisação de guichês —, o modelo de entregas frequentes em sprints quinzenais permite validar hipóteses com operadores reais, enquanto a automação de testes impede regressões transacionais.

\subsecao{Evolução pelas Três Macroentregas}

O ciclo de vida corporativo é governado por três marcos contratuais evolutivos:
\begin{listaetapas}
  \itemetapa{Macroentrega 1 (v1.0 MVP Core Operacional)}{Implementação da fundação do sistema: arquitetura Next.js 16, modelagem relacional Prisma, catálogo de filmes, sessões de cinema, reserva web com bloqueio atômico de 10 minutos (\codigoinline{SeatHold}), venda presencial no guichê e painel administrativo.}
  \itemetapa{Macroentrega 2 (v2.0 Escala e Integrações Financeiras)}{Acoplamento de gateways de pagamento online (PIX dinâmico e Stripe), integração do serviço transacional de mensageria (\textbf{Resend}) para entrega de vouchers com QR Code e suporte a alta concorrência em nuvem.}
  \itemetapa{Macroentrega 3 (v3.0 Multi-Tenant Enterprise)}{Expansão para redes exibidoras com módulos de faturamento automático, emissão fiscal de consumidor (NFC-e), ponto de venda de bomboniere e observabilidade distribuída.}
\end{listaetapas}

% DIAGRAMA DO CICLO DE VIDA DO PROJETO IMPORTADO DE ASSETS/CAPITULO-03/DIAGRAMAS/
\figuracomfonte{capitulo-03/diagramas/diagrama-ciclo-de-vida.pdf}{0.95\textwidth}{Fluxo do Ciclo de Vida do Projeto CineDoom (Sprints Ágeis e Release)}{CineSoft Solutions (2026)}{fig:ciclo_vida}

\secao{Inventário de infraestrutura e ferramentas necessárias}

A infraestrutura técnica foi selecionada para garantir disponibilidade de 99.5\%, segurança em conformidade com a LGPD e alta produtividade da equipe de desenvolvimento.
\begin{center}
\small
\begin{tabularx}{\textwidth}{l p{2.8cm} X}
\toprule
\textbf{Categoria} & \textbf{Ferramenta Adotada} & \textbf{Finalidade no Ciclo de Vida} \\
\midrule
Controle de Versão & GitHub Enterprise & Hospedagem de código, code review e branch protection rules. \\
Esteira de CI/CD & GitHub Actions & Execução automatizada de lint, typecheck, vitest e deploy. \\
Banco de Dados & PostgreSQL 16 (Neon) & Armazenamento relacional com transações ACID estritas. \\
Ambiente Local & Node.js 20 LTS / Docker & Padronização de runtime entre todos os desenvolvedores. \\
\bottomrule
\end{tabularx}
\end{center}

\secao{Justificativa estratégica e viabilidade do projeto}

O mercado cinematográfico brasileiro é composto por mais de 3.400 salas, das quais aproximadamente 35\% pertencem a circuitos independentes ou pequenos exibidores que operam com softwares legados de alto custo de licença ou planilhas despadronizadas.

O CineDoom viabiliza a modernização desses estabelecimentos com custo previsível, eliminando as filas de bilheteria e aumentando a taxa de ocupação das salas em até 22\% por meio de reservas online antecipadas.

\secao{Papéis, competências e matriz RACI}

A conformidade com o nível G do MPS.BR exige a clara segregação de responsabilidades no time de engenharia (\textbf{Equipe Doomscoders}):
{\footnotesize
\begin{tabularx}{\textwidth}{>{\hsize=1.000\hsize\sloppy\arraybackslash}X C{1.69cm} C{1.52cm} C{1.52cm} C{1.69cm}}
\arrayrulecolor{GithubTableBorder}\toprule
\rowcolor{BrandTableHeader}
\textbf{Atividade do Ciclo de Vida} & \textbf{Gabriel Lourenço} & \textbf{Carlos Eduardo} & \textbf{Kayque Vicente} & \textbf{Yury Nogueira} \\
\midrule
\arrayrulecolor{BrandGrayRule}
Planejamento de Sprints, Cronograma e Riscos & C & C & \textbf{A / R} & C \\
\midrule
Elicitação de Requisitos e Especificações Técnicas (SDD) & C & I & C & \textbf{A / R} \\
\midrule
Arquitetura da Solução e Decisões Técnicas (ADRs) & \textbf{A / R} & C & I & C \\
\midrule
Modelagem de Dados e Isolamento de Tenants & \textbf{A / R} & R & I & C \\
\midrule
Criptografia de Dados Pessoais e Conformidade LGPD & C & \textbf{A / R} & C & R \\
\midrule
Implementação do Motor de Reservas e Bilheteria & \textbf{A / R} & R & C & I \\
\midrule
Automação de CI/CD e Infraestrutura de Nuvem & C & \textbf{A / R} & C & I \\
\midrule
Criação e Execução de Testes Automatizados (Vitest) & R & C & \textbf{A / R} & C \\
\midrule
Governança de Qualidade e Revisão de Processos & C & I & \textbf{A / R} & C \\
\midrule
Documentação Técnica e de Processos & C & C & C & \textbf{A / R} \\
\midrule
Gestão de Configuração (Git) e Revisão de Código & \textbf{A} & R & R & I \\
\midrule
Suporte e Manutenção Pós-Implantação & C & \textbf{A / R} & R & I \\
\arrayrulecolor{GithubTableBorder}\bottomrule
\end{tabularx}%
}


\secao{Plano de qualidade com critérios e métricas}

O plano de qualidade da CineSoft Solutions estabelece três pilares fundamentais:
\begin{listapontos}
  \itemponto{Cobertura de Testes}{Nenhum commit é integrado ao branch \codigoinline{main} sem que a suíte automatizada mantenha cobertura superior a 80\% nos módulos de reserva e pagamento.}
  \itemponto{Desempenho de Checkout}{Tempo total de transação de bloqueio de assento inferior a 400 milissegundos na bilheteria presencial.}
  \itemponto{Disponibilidade Contratual}{Taxa de uptime do serviço em nuvem superior a 99.5\% mensurada por monitoramento sintético.}
\end{listapontos}

\secao{Estratégia de gestão do conhecimento}

Todo o conhecimento de engenharia do CineDoom é mantido versionado no próprio repositório na forma de:
\begin{listapontos}
  \itemponto{Architecture Decision Records (ADRs)}{Registro formal em Markdown de todas as escolhas técnicas relevantes (banco, framework, mensageria).}
  \itemponto{Manual de Processos Vivo}{Documentação viva e auditável refletida em sincronia com o código-fonte em produção.}
\end{listapontos}

\verificacaoconformidade{Ponto de Controle — Processos Organizacionais}{%
  \begin{enumerate}[label=\textbf{\alph*)}, leftmargin=0.8cm, itemsep=3pt]
    \item Como a matriz RACI implementada na CineSoft previne gargalos operacionais e conflitos de competência técnica?
    \item De que forma as decisões registradas em ADRs apoiam a conformidade com a cláusula 6.2.6 (\textit{Knowledge Management}) da ISO 12207?
  \end{enumerate}
}

\notatecnica{%
  \itemcaixa{A utilização de TypeScript ponta a ponta (do frontend Next.js ao schema Prisma) reduz em até 40\% os erros de integração entre cliente e servidor, unificando a modelagem de domínio.}
}

\saibamaisnormativo{%
  \itemsaibamais{Consulte as diretrizes do \linkcaixa{https://www.iso.org/standard/63711.html}{ISO/IEC/IEEE 12207:2021 (Cláusula 6.2)} para aprofundar-se nos processos de gestão de portfólio, infraestrutura e recursos humanos.}
}

\sintesecapitulo{%
  \itemcaixa{Processos organizacionais bem estabelecidos são a garantia de que o CineDoom evoluirá de forma consistente, sustentável e com governança auditável.}
}
